\documentclass[10pt,a4paper]{amsart}
\usepackage[margin=19mm]{geometry}
\usepackage{amsmath,amssymb,mathtools}
\usepackage[scr=rsfs,cal=euler]{mathalfa}
\usepackage{xfrac}
\usepackage[unicode,hidelinks,bookmarks=false]{hyperref}
\newcommand{\ie}{i.e.\ }
\newcommand{\cf}{cf.\ }
\newcommand{\resp}{respectively\ }
\newcommand{\Subsection}{\subsection}
\providecommand{\nobreakdash}{\nobreak\hbox{-}}
\setlength{\emergencystretch}{2em}
\multlinegap0pt
\usepackage{bm}
\usepackage{bbm}
\usepackage{dsfont}
\usepackage{xspace}
\usepackage{cancel}
\usepackage{mathtools}
\usepackage{cleveref}
\usepackage{amsthm}
\makeatletter
\@ifundefined{thm}{\newtheorem{thm}{Thm}}{}
\@ifundefined{df}{\newtheorem{df}[thm]{Definition}}{}
\@ifundefined{lem}{\newtheorem{lem}[thm]{Lemma}}{}
\makeatother
\newcommand{\Osp}{\mathcal{OS}}
\title[The Pontryagin-Thom theorem for families of framed equivaria]{The Pontryagin-Thom theorem for families of framed equivariant manifolds}
\author{Lucas Williams}
\date{Source: arXiv:2503.20818v3 (CC BY 4.0)}
\begin{document}
\maketitle

\noindent\emph{Source and licence.} The Pontryagin-Thom theorem for families of framed equivariant manifolds. Version arXiv:2503.20818v3 (CC BY 4.0), distributed under the Creative Commons Attribution 4.0 International licence, as recorded in the arXiv licence metadata for this exact version and shown on the abstract page https://arxiv.org/abs/2503.20818v3. Full rights notice: licenses/arxiv-2503.20818-CC-BY-4.0.txt.\par\smallskip

\noindent\textit{Editorial scope.} Counted unit: the Lem whose source label is \texttt{lem:quotient space} in the article (source0.tex lines 1199--1208, source label \texttt{lem:quotient space}), delivered together with the article's own complete original proof of it (source0.tex lines 1210--1248, one \texttt{proof} environment of 39 lines). No mathematics is written, repaired or re-derived: statement and proof are the article's own text line for line. Required context retained uncounted: df:genuine and categorical fixed points (lines 744--746); df:geometric fixed points (lines 1182--1195). Every definition, display and label of those lines is retained unchanged. Omitted from this package and not redistributed: 1--743, 747--1181, 1196--1198, 1209--1209, 1249--2393. Nothing inside a retained excerpt is omitted or altered: every retained body is delivered exactly as the article's lines are written, so no omission receipt of any kind accompanies this package. Citations: the retained excerpts carry no citation command at all, so no bibliography entry of the article is transcribed and no reference is redistributed. Reference adaptation: none was needed and none was made; every reference inside the delivered unit resolves inside it, so no unresolved reference or citation is produced. Printed slips kept literal: no printed word, symbol, hypothesis or number of the counted statement or of its proof was altered. Layout and class: the article's own class is \texttt{amsart} with packages including amsmath, amssymb, amsthm, array, color, multirow, pdflscape, mathrsfs, subcaption, todonotes, tikz, tikz-cd, float, imakeidx; the wrapper is the lane's pinned \texttt{amsart} editorial wrapper at 10pt on A4, which restates the theorem-like environments the retained text needs and adds the article's own macro definitions that the retained text needs (\texttt{\textbackslash Osp}), with extra wrapper packages (bm, bbm, dsfont, xspace, cancel, mathtools, cleveref). The delivered numbering is the unit's own. Assets: no retained span contains a figure environment, a graphics-inclusion command, a PDF-page inclusion, a table or a TikZ picture; no image file is copied into this package, the package declares no asset dependency and no image file is redistributed with it. Licence: version arXiv:2503.20818v3 of this article is distributed under the Creative Commons Attribution 4.0 International licence, as recorded in the arXiv licence metadata for this exact version and shown on the abstract page https://arxiv.org/abs/2503.20818v3. A full rights notice is delivered at licenses/arxiv-2503.20818-CC-BY-4.0.txt. Characters: every retained excerpt is pure ASCII, so the delivered engine, XeLaTeX with its default Latin Modern text font, reports no missing character.
\begin{df}\label{df:genuine and categorical fixed points}
Let $X\in G\Osp$. The \textbf{categorical fixed points} of $X$ is the orthogonal $WH$-spectrum with $n$th level given by $X_n^H$. The \textbf{genuine fixed points} of $X$, denoted $X^G$, are the right derived categorical fixed points. This means that $X^G$ is constructed by replacing $X$ with a stably equivalent fibrant spectrum and then taking the categorical fixed points.
\end{df}

\begin{df}\label{df:geometric fixed points}
Let $H\leq G$ and let $X$ be a cofibrant genuine $G$-spectrum. Define the $G$-space $\widetilde{EP_H}$ up to homotopy as any $G$-CW complex with fixed points of the following type:
\[
\widetilde{EP_H}^K \simeq\begin{cases}
S^0 & \text{ if } H\leq K\\
* & \text{ else. }
\end{cases}
\]
Then the \textbf{geometric fixed points of $X$} are 
\[
(\widetilde{EP_H} \wedge X)^{H}.
\]
Here $\wedge$ denotes smashing $\widetilde{EP_H}$ with every spectrum level of $X$, and $(-)^H$ denotes the genuine fixed points of \cref{df:genuine and categorical fixed points}. 
\end{df}

\begin{lem}\label{lem:quotient space}
Let $F$ and $F'$ be two families of subgroups of $G$ which differ by the conjugacy class of $(H)$ i.e.
\[
F= F'\cup (H)
\]
There is a weak equivalence of $G$-spaces
\[
EF/EF' \simeq G\wedge_{NH} EWH_+ \wedge \widetilde{EP_H}
\]
\end{lem}

\begin{proof}
First recall that the equivariant homotopy type of $EF$ is defined by 
\[
(EF)^K = \begin{cases}
* & \text{ if } K\in F\\
\emptyset & \text{ else }
\end{cases}
\]
Thus, the equivariant homotopy type of $EF/EF'$ is given by 
\[
(EF/EF')^K = \begin{cases}
S^0 & \text{ if } K = H\\
* & \text{ else }
\end{cases} 
\]

Note that the isotropy group of any non-basepoint element of $G\wedge_{NH}EWH_+$ is conjugate to $H$. Moreover, 
\[
(G\wedge_{NH}EWH_+)^H \cong NH\wedge_{NH}EWH_+ \simeq S^0
\]
where the last map is a weak equivalence of non-equivariant spaces. In particular, we have now obtained
\[
(G\wedge_{NH}EWH_+)^K = \begin{cases}
S^0 & \text{ if } K=H\\
* & \text{ if $K$ properly contains $H$}
\end{cases}
\]
combining this with the definition of $\widetilde{EP}_H$ from \cref{df:geometric fixed points} we obtain
\[
(G\wedge_{NH}EWH_+ \wedge\widetilde{EP}_H)^K = \begin{cases}
S^0 & \text{ if } K = H\\
* & \text{ else }
\end{cases}
\]
Therefore, by Elmendorf's theorem, there is a weak equivalence of $G$-spaces
\[
G\wedge_{NH}EWH_+ \wedge\widetilde{EP}_H \simeq EF/EF'
\]
\end{proof}

\end{document}
